\documentclass[11pt,a4paper]{article}

\usepackage[utf8]{inputenc}
\usepackage[T1]{fontenc}
\usepackage[ngerman]{babel}

\usepackage{amsmath,amssymb,amsthm,mathtools}
\usepackage[a4paper,margin=2.4cm]{geometry}
\usepackage{lmodern}
\usepackage{microtype}
\usepackage{hyperref}
\usepackage{enumitem}
\usepackage{booktabs}

\title{Das EABC-Gitter modulo \(12\):\\
	Zwillinge, Drillinge und blockübergreifende Vierlinge}
\author{Kollaborative Ausarbeitung}
\date{\today}

\newtheorem{satz}{Satz}
\newtheorem{lemma}[satz]{Lemma}
\newtheorem{korollar}[satz]{Korollar}
\newtheorem{proposition}[satz]{Proposition}
\theoremstyle{definition}
\newtheorem{definition}[satz]{Definition}
\newtheorem{beispiel}[satz]{Beispiel}
\theoremstyle{remark}
\newtheorem{bemerkung}[satz]{Bemerkung}

\newcommand{\N}{\mathbb{N}}
\newcommand{\Z}{\mathbb{Z}}
\newcommand{\pp}{\mathbb{P}}
\newcommand{\ind}{\mathbf{1}}

\begin{document}
	\maketitle
	
	\begin{abstract}
		Wir betrachten die Zerlegung der natürlichen Zahlen \(>3\) nach ihren Restklassen modulo \(12\)
		in die vier Familien \(E,A,B,C\). Auf dieser Basis wird das zugehörige EABC-Gitter
		\[
		Q_n=(12n-11,\,12n-7,\,12n-5,\,12n-1)
		\]
		eingeführt. Es wird gezeigt, dass alle Primzahlzwillinge \(>3\) eindeutig in genau zwei
		Kanälen des Gitters liegen, dass alle Primzahldrillinge vollständig durch die beiden
		\(3\)-verankerten Ausnahmefälle \((3,5,7)\) und \((3,7,11)\) beschrieben sind, und dass die
		gleichzeitige Aktivierung beider Zwillingskanäle genau den blockübergreifenden
		Primzahlvierlingen entspricht. Damit erhält man eine vollständige lokale Strukturtheorie
		dieser niedrigen Primzahlkonstellationen im modulo-\(12\)-Gitter.
	\end{abstract}
	
	\section{Einleitung}
	
	Jede Primzahl \(p>3\) ist weder durch \(2\) noch durch \(3\) teilbar. Daher kann sie modulo \(12\)
	nur in einer der vier Restklassen
	\[
	1,\ 5,\ 7,\ 11
	\]
	liegen. Diese vier Klassen definieren ein natürliches vierteiliges Gitter der Primzahlen.
	Schreibt man sie in der Reihenfolge
	\[
	E \equiv 1,\qquad A \equiv 5,\qquad B \equiv 7,\qquad C \equiv 11 \pmod{12},
	\]
	so entstehen lokale Viererblöcke
	\[
	Q_n=(E_n,A_n,B_n,C_n),
	\]
	innerhalb derer sich die niedrigsten Primzahlkonstellationen --- insbesondere Zwillinge,
	Drillinge und Vierlinge --- in bemerkenswert straffer Weise organisieren lassen.
	
	Ziel dieses Textes ist eine konsistente Definition--Lemma--Satz--Beweis-Folge, in der diese
	lokale Struktur vollständig beschrieben wird.
	
	\section{Die vier Familien}
	
	\begin{definition}[EABC-Familien]
		Für \(n\in\N\) definieren wir
		\[
		E_n:=12n-11,\qquad
		A_n:=12n-7,\qquad
		B_n:=12n-5,\qquad
		C_n:=12n-1.
		\]
		Wir nennen
		\[
		Q_n:=(E_n,A_n,B_n,C_n)
		\]
		den \(n\)-ten \emph{EABC-Vierertruppel}.
	\end{definition}
	
	\begin{lemma}[Vollständige Zerlegung der Primzahlen \(>3\)]
		Jede Primzahl \(p>3\) liegt eindeutig in genau einer der vier Familien
		\[
		E,\quad A,\quad B,\quad C.
		\]
	\end{lemma}
	
	\begin{proof}
		Ist \(p>3\) prim, so ist \(p\) weder durch \(2\) noch durch \(3\) teilbar. Daher kann
		\[
		p \equiv 1,5,7,\text{ oder }11 \pmod{12}
		\]
		gelten und keine andere Restklasse. Die vier Folgen
		\[
		12n-11,\quad 12n-7,\quad 12n-5,\quad 12n-1
		\]
		parametrisieren diese vier Klassen disjunkt und vollständig. Also liegt jede Primzahl \(p>3\)
		eindeutig in genau einer dieser Familien.
	\end{proof}
	
	\begin{korollar}
		Die Menge der Primzahlen \(>3\) zerfällt disjunkt in
		\[
		\pp\setminus\{2,3\}
		=
		(\pp\cap E)\,\dot\cup\,(\pp\cap A)\,\dot\cup\,(\pp\cap B)\,\dot\cup\,(\pp\cap C).
		\]
	\end{korollar}
	
	\section{Die beiden Zwillingskanäle}
	
	\begin{definition}[Zwillingskanäle]
		Zu jedem Vierertruppel
		\[
		Q_n=(E_n,A_n,B_n,C_n)
		\]
		definieren wir zwei natürliche Zwillingskanäle:
		\[
		T_n^{\mathrm{int}}:=(A_n,B_n)=(12n-7,\,12n-5),
		\]
		\[
		T_n^{\mathrm{rand}}:=(C_n,E_{n+1})=(12n-1,\,12n+1).
		\]
		Weiter definieren wir die \emph{Zwillingssignatur}
		\[
		\tau(Q_n)
		:=
		\bigl(
		\ind_{T_n^{\mathrm{int}}\text{ ist Primzahlzwilling}},
		\ind_{T_n^{\mathrm{rand}}\text{ ist Primzahlzwilling}}
		\bigr)
		\in\{0,1\}^2.
		\]
	\end{definition}
	
	\begin{lemma}[Es gibt genau zwei Zwillingskanäle]
		Jeder Primzahlzwilling \((p,p+2)\) mit \(p>3\) liegt eindeutig in genau einem der beiden Kanäle:
		\[
		(p,p+2)=(A_n,B_n)
		\quad\text{oder}\quad
		(p,p+2)=(C_n,E_{n+1})
		\]
		für ein eindeutig bestimmtes \(n\).
	\end{lemma}
	
	\begin{proof}
		Sei \((p,p+2)\) ein Primzahlzwilling mit \(p>3\). Da beide Zahlen prim sind, sind sie ungerade
		und nicht durch \(3\) teilbar. Modulo \(12\) kommen daher nur die Muster
		\[
		(5,7)\quad\text{oder}\quad(11,1)
		\]
		in Frage.
		
		Im ersten Fall gilt
		\[
		p=12n-7=A_n,\qquad p+2=12n-5=B_n
		\]
		für ein eindeutig bestimmtes \(n\).
		
		Im zweiten Fall gilt
		\[
		p=12n-1=C_n,\qquad p+2=12n+1=E_{n+1}
		\]
		für ein eindeutig bestimmtes \(n\).
		
		Andere Restklassenmuster sind für Primzahlzwillinge \(>3\) unmöglich.
	\end{proof}
	
	\begin{satz}[Vollständige Klassifikation der Primzahlzwillinge \(>3\)]
		Die Primzahlzwillinge \(>3\) sind genau die primen Realisierungen der beiden linearen Muster
		\[
		(12n-7,12n-5)
		\qquad\text{und}\qquad
		(12n-1,12n+1).
		\]
	\end{satz}
	
	\begin{proof}
		Dies folgt unmittelbar aus dem vorigen Lemma.
	\end{proof}
	
	\begin{bemerkung}
		Damit ist jeder Primzahlzwilling \(>3\) im EABC-Gitter entweder ein
		\emph{interner Zwilling} vom Typ \(A\!-\!B\) oder ein \emph{Randzwilling} vom Typ \(C\!-\!E\).
	\end{bemerkung}
	
	\section{Primzahldrillinge}
	
	\begin{definition}[Dichte Drillingsmuster]
		Unter einem \emph{Primzahldrilling} verstehen wir ein Primzahltripel einer der beiden Formen
		\[
		(p,p+2,p+6)
		\qquad\text{oder}\qquad
		(p,p+4,p+6).
		\]
	\end{definition}
	
	\begin{lemma}[Jeder Primzahldrilling ist \(3\)-verankert]
		Jeder Primzahldrilling enthält notwendig die Primzahl \(3\).
	\end{lemma}
	
	\begin{proof}
		Betrachte eines der beiden Muster
		\[
		p,\ p+2,\ p+6
		\qquad\text{oder}\qquad
		p,\ p+4,\ p+6.
		\]
		Modulo \(3\) enthält ein solches Tripel stets mindestens eine Zahl, die zu \(0\pmod 3\)
		kongruent ist. Soll diese Zahl prim sein, so muss sie gleich \(3\) sein.
	\end{proof}
	
	\begin{satz}[Vollständige Klassifikation der Primzahldrillinge]
		Die einzigen Primzahldrillinge sind
		\[
		(3,5,7)
		\qquad\text{und}\qquad
		(3,7,11).
		\]
	\end{satz}
	
	\begin{proof}
		Nach dem vorherigen Lemma muss jeder Primzahldrilling die Primzahl \(3\) enthalten.
		Dann bleiben für die beiden dichten Muster nur die beiden Möglichkeiten
		\[
		(3,5,7)
		\qquad\text{und}\qquad
		(3,7,11).
		\]
		Beide Tripel bestehen tatsächlich aus Primzahlen. Weitere Fälle gibt es nicht.
	\end{proof}
	
	\begin{korollar}
		Nach Abspaltung des singulären Ankers \(3\) entsprechen die beiden Primzahldrillinge genau den
		Familienpaaren
		\[
		(A_1,B_1)=(5,7)
		\qquad\text{und}\qquad
		(B_1,C_1)=(7,11).
		\]
	\end{korollar}
	
	\begin{bemerkung}
		Primzahldrillinge sind damit keine freien inneren Konfigurationen des EABC-Gitters, sondern
		vollständig durch den Ausnahmepunkt \(3\) erzwungene Randphänomene.
	\end{bemerkung}
	
	\section{Doppelstellen und blockübergreifende Vierlinge}
	
	\begin{definition}[Doppelstelle]
		Ein Vierertruppel \(Q_n\) heißt \emph{Doppelstelle}, wenn beide Zwillingskanäle aktiv sind, also
		\[
		\tau(Q_n)=(1,1).
		\]
		Das ist äquivalent dazu, dass
		\[
		A_n,\ B_n,\ C_n,\ E_{n+1}
		\]
		alle prim sind.
	\end{definition}
	
	\begin{lemma}[Doppelstellen sind genau blockübergreifende Vierlinge]
		Ein Vierertruppel \(Q_n\) ist genau dann eine Doppelstelle, wenn
		\[
		(12n-7,\ 12n-5,\ 12n-1,\ 12n+1)
		\]
		ein Primzahlvierling ist.
	\end{lemma}
	
	\begin{proof}
		Nach Definition bedeutet
		\[
		\tau(Q_n)=(1,1),
		\]
		dass sowohl
		\[
		(12n-7,12n-5)
		\]
		als auch
		\[
		(12n-1,12n+1)
		\]
		Primzahlzwillinge sind. Dies ist genau dann der Fall, wenn alle vier Zahlen
		\[
		12n-7,\ 12n-5,\ 12n-1,\ 12n+1
		\]
		prim sind. Diese bilden das Muster
		\[
		p,\ p+2,\ p+6,\ p+8,
		\]
		also einen Primzahlvierling.
	\end{proof}
	
	\begin{satz}[Vierlinge als Kopplung beider Zwillingskanäle]
		Die Doppelstellen des EABC-Gitters sind genau die blockübergreifenden Primzahlvierlinge.
		Insbesondere ist ein Primzahlvierling im EABC-Gitter genau die lokale Konfiguration, in der
		\[
		(A_n,B_n)
		\quad\text{und}\quad
		(C_n,E_{n+1})
		\]
		gleichzeitig als Primzahlzwillinge auftreten.
	\end{satz}
	
	\begin{proof}
		Dies folgt unmittelbar aus dem vorherigen Lemma.
	\end{proof}
	
	\section{Hauptsatz}
	
	\begin{satz}[Lokale Vollständigkeit des EABC-Gitters]
		Das EABC-Gitter
		\[
		Q_n=(12n-11,12n-7,12n-5,12n-1)
		\]
		liefert eine vollständige lokale Strukturtheorie der Primzahlkonstellationen niedriger Ordnung:
		
		\begin{enumerate}[label=\textup{(\roman*)}]
			\item Alle Primzahlen \(>3\) liegen eindeutig in genau einer der vier Familien
			\[
			E,\ A,\ B,\ C.
			\]
			
			\item Alle Primzahlzwillinge \(>3\) entstehen eindeutig in genau einem von zwei Kanälen:
			\[
			(A_n,B_n)=(12n-7,12n-5)
			\quad\text{oder}\quad
			(C_n,E_{n+1})=(12n-1,12n+1).
			\]
			
			\item Alle Primzahldrillinge sind singuläre \(3\)-verankerte Randphänomene und bestehen genau aus
			\[
			(3,5,7)
			\quad\text{und}\quad
			(3,7,11).
			\]
			
			\item Alle lokalen Doppelkopplungen beider Zwillingskanäle sind genau die blockübergreifenden
			Primzahlvierlinge
			\[
			(12n-7,12n-5,12n-1,12n+1).
			\]
		\end{enumerate}
		
		Damit ist die gesamte lokale Zwillings-, Drillings- und Vierlingskopplungsstruktur des
		Primzahlgitters modulo \(12\) vollständig beschrieben.
	\end{satz}
	
	\begin{proof}
		Aussage \textup{(i)} ist Lemma 1. Aussage \textup{(ii)} ist der Klassifikationssatz über die
		beiden Zwillingskanäle. Aussage \textup{(iii)} ist der Klassifikationssatz der Primzahldrillinge.
		Aussage \textup{(iv)} ist der Vierlingssatz über die Doppelstellen. Damit sind alle genannten
		lokalen Konfigurationsklassen vollständig erfasst.
	\end{proof}
	
	\section{Beispiele}
	
	\begin{beispiel}[Zwilling im inneren Kanal]
		Für \(n=2\) ist
		\[
		Q_2=(13,17,19,23).
		\]
		Hier ist
		\[
		(A_2,B_2)=(17,19)
		\]
		ein Primzahlzwilling.
	\end{beispiel}
	
	\begin{beispiel}[Zwilling am Randkanal]
		Für \(n=5\) ist
		\[
		Q_5=(49,53,55,59),
		\]
		und
		\[
		(C_5,E_6)=(59,61)
		\]
		ist ein Primzahlzwilling.
	\end{beispiel}
	
	\begin{beispiel}[Drilling]
		Der Primzahldrilling
		\[
		(3,5,7)
		\]
		entspricht nach Abspaltung von \(3\) genau dem inneren Zwillingspaar
		\[
		(A_1,B_1)=(5,7).
		\]
	\end{beispiel}
	
	\begin{beispiel}[Doppelstelle]
		Für \(n=9\) gilt
		\[
		Q_9=(97,101,103,107).
		\]
		Dann sind
		\[
		(A_9,B_9)=(101,103)
		\]
		und
		\[
		(C_9,E_{10})=(107,109)
		\]
		beide Primzahlzwillinge. Also ist \(Q_9\) eine Doppelstelle, und
		\[
		(101,103,107,109)
		\]
		ist ein blockübergreifender Primzahlvierling.
	\end{beispiel}
	
	\section{Schlussbemerkung}
	
	Die modulo-\(12\)-Zerlegung in die vier Familien \(E,A,B,C\) erzeugt ein natürliches lokales Gitter,
	in dem die kleinsten Primzahlkonstellationen eine bemerkenswert einfache und vollständige Struktur
	besitzen:
	
	\begin{itemize}[itemsep=0.25em]
		\item Primzahlen \(>3\) füllen genau die vier Familien,
		\item Primzahlzwillinge laufen genau über zwei Kanäle,
		\item Primzahldrillinge sind vollständig durch den singulären Anker \(3\) bestimmt,
		\item und Primzahlvierlinge erscheinen genau als Doppelkopplung beider Zwillingskanäle.
	\end{itemize}
	
	Damit wird das EABC-Gitter zu einer sehr kompakten lokalen Beschreibung der niedrigsten
	Primzahlmuster.
	
\end{document}